\uuid{WY6r}
\exo7id{3712}
\titre{exo7 3712}
\auteur{quercia}
\organisation{exo7}
\datecreate{2010-03-11}
\isIndication{false}
\isCorrection{true}
\chapitre{Espace euclidien, espace normé}
\sousChapitre{Espace vectoriel euclidien de dimension 3}
\cpge{mpsi-mp2i.matrices-applications-lineaires.a, pcsi.matrices-applications-lineaires.a, ptsi.matrices-applications-lineaires.a, mp-mpi.endomorphismes-euclidien.b, pc.endomorphismes-euclidiens.b, psi.endomorphismes-euclidien.b, pt.isometries-euclidien.b, mp-mpi.endomorphismes-euclidien.e, psi.endomorphismes-euclidien.e, pt.isometries-euclidien.d}
\module{Algèbre}
\niveau{L2}
\difficulte{}

\contenu{
\question{
Déterminer la matrice de la rotation $\cal R$ de $\R^3$ dans une base
orthonormée $(\vec i,\vec j,\vec k\,)$ telle que ${\cal R}(\vec u)=\vec u$
avec $\vec u\bigl(\frac1{\sqrt3},\frac{-1}{\sqrt3},\frac1{\sqrt3}\bigr)$
et ${\cal R}(\vec i) = \vec k$. Donner son angle de rotation.
}
\reponse{
$M=\begin{pmatrix}0 &-1 &0\cr 0 &0 &-1\cr 1 &0 &0\cr\end{pmatrix}$, $\theta = \frac{2\pi}3$.
}
}
